\documentclass{article}
\usepackage{amsmath, amssymb, booktabs}

\begin{document}

\title{Practice Questions on the Ricardian Model}
\author{}
\date{}
\maketitle

\section*{Problem 1: Comparative Advantage and Opportunity Cost}
Two countries, \textbf{A} and \textbf{B}, produce two goods: \textbf{Wheat} and \textbf{Cloth}. The labor productivity (output per worker per hour) is given below:

\begin{table}[h]
    \centering
    \begin{tabular}{lcc}
        \toprule
        Country & Wheat (units per hour) & Cloth (units per hour) \\
        \midrule
        \textbf{A} & 10 & 5 \\
        \textbf{B} & 4 & 2 \\
        \bottomrule
    \end{tabular}
\end{table}

\begin{enumerate}
    \item[(a)] Calculate the opportunity cost of producing one unit of \textbf{wheat} in both countries.
    \item[(b)] Calculate the opportunity cost of producing one unit of \textbf{cloth} in both countries.
    \item[(c)] Which country has a \textbf{comparative advantage} in wheat production? Which has a comparative advantage in cloth production?
\end{enumerate}

\section*{Problem 2: Gains from Specialization and Trade}
Consider two countries, \textbf{X} and \textbf{Y}, producing \textbf{Cars} and \textbf{Computers}. The total labor force is \textbf{1,000 workers} in each country. The labor productivity is:

\begin{table}[h]
    \centering
    \begin{tabular}{lcc}
        \toprule
        Country & Cars per worker per year & Computers per worker per year \\
        \midrule
        \textbf{X} & 4 & 8 \\
        \textbf{Y} & 2 & 6 \\
        \bottomrule
    \end{tabular}
\end{table}

\begin{enumerate}
    \item[(a)] Calculate the opportunity cost of producing \textbf{one car} in terms of computers in each country.
    \item[(b)] If each country fully specializes in the good in which it has a comparative advantage, how many total \textbf{cars} and \textbf{computers} will be produced?
    \item[(c)] Show that trade allows each country to consume beyond its production possibilities frontier (PPF).
\end{enumerate}

\section*{Problem 3: Terms of Trade and Trade Gains}
Two countries, \textbf{Brazil} and \textbf{Canada}, produce \textbf{Coffee} and \textbf{Timber}. The labor productivity is as follows:

\begin{table}[h]
    \centering
    \begin{tabular}{lcc}
        \toprule
        Country & Coffee (tons per worker) & Timber (tons per worker) \\
        \midrule
        \textbf{Brazil} & 12 & 4 \\
        \textbf{Canada} & 8 & 6 \\
        \bottomrule
    \end{tabular}
\end{table}

\begin{enumerate}
    \item[(a)] Find the opportunity cost of producing \textbf{1 ton of coffee} in each country.
    \item[(b)] Suppose Brazil specializes in \textbf{coffee} and Canada specializes in \textbf{timber}. If they trade at a \textbf{terms of trade} of 1 ton of coffee = \textbf{0.5 tons of timber}, will both countries benefit? Explain why or why not.
\end{enumerate}

\section*{Problem 4: Effect of Technological Change on Trade Patterns}
Initially, two countries, \textbf{India} and \textbf{Germany}, produce \textbf{Steel} and \textbf{Textiles}. The productivity levels are:

\begin{table}[h]
    \centering
    \begin{tabular}{lcc}
        \toprule
        Country & Steel (tons per worker) & Textiles (meters per worker) \\
        \midrule
        \textbf{India} & 5 & 20 \\
        \textbf{Germany} & 10 & 30 \\
        \bottomrule
    \end{tabular}
\end{table}

\begin{enumerate}
    \item[(a)] Determine each country's comparative advantage before any technological change.
    \item[(b)] Suppose a technological improvement in Germany \textbf{doubles its steel productivity} but does not affect textiles. How does this affect Germany’s comparative advantage?
    \item[(c)] How would this change affect the \textbf{global pattern of trade} between India and Germany?
\end{enumerate}

\end{document}
